%% ============================================================
%% chapters/introduction.tex
%% Introduction — Vedant Patil Master's Thesis
%% ============================================================

\chapter{Introduction}
\label{ch:introduction}

\noindent Artificial intelligence is leaving the server room.
 
For most of its recent history, the infrastructure story of machine learning was simple: train on a cluster, serve from a cloud, assume a network. That story worked well enough while AI lived inside data centers and the people who needed its outputs could wait seconds or minutes for a round-trip to a backend. But AI is increasingly being asked to operate somewhere else entirely. On a hillside. In a field. At a remote sensing station hours from the nearest cell tower. In environments where the assumption of connectivity is not just occasionally wrong but structurally, permanently wrong.
 
The shift matters because the tools that built modern AI systems were not designed for this. Container orchestration frameworks were designed for horizontally scalable cloud services, not for a single laptop deployed in a farm field. Package registries, image caches, and DNS resolution were designed under the assumption that a network path to the outside world exists. Even the debugging workflows used during development, running a system locally and then assuming it will behave the same way elsewhere, carry implicit assumptions about environmental consistency that physical deployment sites simply do not provide.
 
What happens when those assumptions meet reality is not a software bug. The code may be entirely correct. The model may be well-trained and well-tested. What fails is the layer underneath: the infrastructure contracts the application was written against, which the deployment environment does not honor. The system does not gradually degrade. It collapses at initialization, before a single inference is ever attempted, and it does so in ways that are invisible during development because development environments are, by design, nothing like the field.
 
This is a systems problem, not an algorithms problem. And it is a solved problem only if you treat it as one.
 
The dominant response in the research community has been to optimize the AI itself for constrained hardware: smaller models, quantized weights, pruned architectures. These are genuine contributions. But they address a different failure mode. A model that runs efficiently on a microcontroller is still useless if the container it lives in cannot start because the orchestrator expected a stable network interface that the deployment site does not have. Efficiency research assumes the platform is working. In far-edge environments, that assumption needs to be earned, not inherited.
 
This thesis is about earning it. Specifically, it is about understanding why far-edge AI systems fail at deployment, classifying those failures in a way that makes them tractable, and designing the architectural principles that prevent them. The argument is not that any particular application domain demands this. It is that any AI system operating at the far edge, regardless of what it is for, faces the same class of infrastructure problem. The solutions are likewise general.
%% ============================================================
\section{Motivation: The Gap Between the Lab and the Field}
\label{sec:intro_motivation}
%% ============================================================
Machine learning models are tested on curated benchmarks, trained on GPU clusters, and validated in controlled environments where infrastructure is predictable and internet connectivity is treated as a given. None of that infrastructure exists at the far edge.

Far-edge deployments operate at the physical location of a sensing event, far removed from cloud infrastructure and frequently from any network at all. These systems cannot offload computation to the cloud because there is no path to the cloud. They cannot recover a failed container by pulling from an image registry because the registry is unreachable. They run on hardware that was never in scope when the tools used to build them were designed.
The result is a deployment gap. Capable AI systems, whose models are correct and whose code passes every test, arrive in the field and fail before inference is ever attempted. Not because of algorithmic weakness. Because the platform underneath them was built on assumptions the field does not honor.

This thesis is about that gap: what produces it, why it is systematic rather than accidental, and how it can be closed through deliberate design.
%% ============================================================
\section{Far-Edge AI as a Distinct Systems Problem}
\label{sec:far_edge_ai}
%% ============================================================
The terms "edge computing" and "far-edge computing" are often used interchangeably in the literature, but the distinction carries real consequences for system design. Near-edge infrastructure includes cellular base stations, regional server rooms, and enterprise gateways. It is reasonably well-provisioned, runs standard software stacks, and maintains reliable paths to cloud backends. Far-edge infrastructure lies beyond that boundary. It is physically remote, power-constrained, and frequently disconnected. Hardware is chosen for weight or power budget, not compute headroom. The people operating it are domain specialists, not platform engineers.

Consider what this means in practice. An aerial imaging system deployed at a remote agricultural site must complete its analysis during the mission window, not after a round-trip to a cloud backend that may not be reachable. An autonomous wildlife monitoring system triggered by a detection event must dispatch within seconds, not the minutes that a human operator or a cloud relay would require. In both cases, the operational constraint is not primarily about model accuracy or inference speed. It is about whether the platform can function at all under the conditions the site imposes.
The systems literature has largely not engaged with this. Most edge AI research focuses on making inference faster or smaller on constrained hardware. That work is valuable. But it assumes the platform under the model is functional. This thesis starts one layer below that assumption.
%% ============================================================
\section{Two Deployments, One Argument}
\label{sec:two_deploy_one_arg}
%% ============================================================
This thesis is grounded in two real field deployments, not laboratory simulations. The first, a containerized agricultural AI system called OpenPASS, was carried to a soybean farm in Stoutsville, Ohio. It failed in ways that were invisible during development. Those failures were analyzed, classified, and resolved through architectural changes. The second, a wildlife monitoring system called the Smart Backpack, was deployed at The Wilds Conservation Center in southeastern Ohio. It was built with the lessons of the first deployment in mind, on different hardware, using a different orchestration framework, in a completely different application domain. It worked.

The relationship between the two deployments is not coincidental. Chapter 3 documents what broke and why. Chapter 4 shows the same design principles applied proactively, before anything broke. The consistency of the outcome across two domains, two hardware platforms, and two orchestration frameworks is the thesis's central empirical claim.
%% ============================================================
\section{Thesis Claim}
\label{sec:thesis_claim}
%% ============================================================
Far-edge AI systems fail at deployment not due to algorithmic weakness, but due to infrastructure heterogeneity. These failures are systematic, classifiable, and addressable through offline-first design principles. \textbf{We define offline-first design as an approach in which external network connectivity is treated as unavailable from the beginning of system design, rather than as a resource to be degraded gracefully from when lost.} Our results are consistent with these principles generalizing across orchestration frameworks, hardware architectures, and application domains.

This claim is validated through two field deployments.

The first is OpenPASS, an agricultural AI system carried to a soybean farm in Stoutsville, Ohio. The initial deployment failed. The failures were analyzed, classified as inflection points, and resolved through architectural changes. The result is a validated offline-first deployment pattern for far-edge K3S applications.

The second is the Smart Backpack system, deployed at The Wilds Conservation Center in southeastern Ohio. This system applied offline-first principles from the beginning of design, before the first line of code was written, on a Raspberry Pi 4B using Docker Compose rather than Kubernetes. Over a 20-hour field deployment it completed 38 autonomous drone missions triggered by camera trap events, with no human involvement in dispatch and no external network dependency during operation.

Same principles. Different stack. Different domain. Results consistent with the principles generalizing.
%% ============================================================
\section{Contributions}
\label{sec:contribution}
%% ============================================================
Inflection-point classification framework.
We introduce the inflection point as a formal concept for far-edge deployment analysis: a site-induced, platform-level failure mode that is distinct from application bugs and requires an architectural response. Three inflection points from the OpenPASS agricultural deployment are characterized, each with its triggering condition, failure mode, and architectural resolution documented.

\textbf{Offline-first K3S architecture.}
We designed and implemented an offline-first deployment architecture for far-edge K3S applications. Three components address the three inflection points: a VXLAN-based overlay that makes pod networking independent of physical network interfaces at a site; pod scheduling with Pod Disruption Budgets that protects inference containers from eviction under resource pressure; and an extended containerd configuration with on-node Helm repositories that reconstructs the full application stack without any external registry access. This architecture was validated at a real agricultural field site.

\textbf{OpenPassLite}.
We designed and implemented a modular drone mission framework built on the Parrot Olympe SDK. Its architecture centers on drop-in mission modules: each flight behavior is a self-contained Python file that the framework loads without modification to its core. Missions for autonomous GPS navigation, orthomosaic survey capture, and return-to-home are implemented and validated in two separate field deployments across two application domains.

\textbf{SmartFields}.
We designed and implemented a fault-tolerant pipeline orchestrator that coordinates the full detection-to-documentation pipeline in the wildlife monitoring deployment. SmartFields receives camera trap events via MQTT, resolves GPS coordinates, dispatches a three-stage drone mission through OpenPassLite, monitors stage completion through log inspection, and enforces fault-tolerance rules that prioritize drone safety when a stage fails. A threading lock enforces single-mission execution at all times.
Proposed inflection-point database.

We propose a structured community repository for documenting far-edge deployment failure modes, their causes, and their architectural resolutions. This artifact does not exist yet; we describe its design and make the case for it as a community infrastructure investment.
Cross-domain validation.

Offline-first design principles were independently applied in two field environments using two orchestration frameworks, two hardware architectures, and two application domains. The 38-mission, 92\% success-rate result at The Wilds, achieved with zero network dependency and zero human intervention in dispatch, provides the cross-domain evidence that these principles are not K3S-specific or agriculture-specific. They are a generalizable design approach.
%% ============================================================
\section{Thesis Organization}
\label{sec:thesis_org}
%% ============================================================
Chapter 2 surveys related work across edge computing systems, container orchestration for constrained hardware, digital agriculture, and wildlife monitoring. It locates this thesis within the existing literature and identifies the gaps that motivate the work.

Chapter 3 presents the OpenPASS agricultural system. It describes the Stoutsville, Ohio deployment context, develops the inflection-point classification framework, presents the three inflection points and their architectural resolutions, and reports the results of the offline-first K3S deployment.

Chapter 4 presents the Smart Backpack wildlife monitoring system. It describes the SmartWilds project context, the eight-container Docker Compose architecture, the design of OpenPassLite and SmartFields, and the field validation results from The Wilds Conservation Center.

Chapter 5 draws both chapters into a unified analysis. It traces how the design principles developed in Chapter 3 appear in the architecture of Chapter 4, examines what that cross-domain consistency implies, and presents the proposed inflection-point database.

Chapter 6 states contributions, discusses limitations, and identifies open problems.